\chapter{Petroleum}
\label{ch:petroleum}

A complete oil chain: find it, drill it, refine it, store it, burn it, and sell it at a forecourt.

\begin{iefork}
Entirely this fork, and the largest single addition in it.
\end{iefork}

\section{Prospecting}

Oil sits in \textbf{reservoirs} beneath the world, which have to be found before they can be
drilled. The \textbf{Seismic Survey Kit} reports bearings and ranges rather than handing you a
coordinate.

\section{Drilling}

A \textbf{Drilling Derrick} sunk over a reservoir brings crude to the surface through a
\textbf{Wellhead}, which is what the derrick leaves behind when it is done. A reservoir is finite:
its flow decays as it empties, and it can be blown out by mishandling --- fit a \textbf{Blowout
Preventer} to the well before that becomes your problem.

\subsection*{Pumpjack}

A $3\times6\times5$ machine built \emph{around} an existing Wellhead --- the well bay is drawn in
its plan but is not a part you supply, and forming the structure will not eat the well. It holds no
fluid of its own; it simply drives the well beneath it, for \textbf{512\,IF/t}.

\begin{ietip}
A pumpjack on a healthy field does nothing, and knows it. While a deposit is still above its
free-flow threshold it already produces at full rate unaided, so the machine idles rather than
charge you for nothing. It starts working when the field starts tiring. In city mode it always
runs, because there reservoirs never deplete.
\end{ietip}

\subsection*{Re-injection Well}

A single block that puts water or gas back downhole and gets a second tranche out of a tiring
field --- the thinking player's alternative to burning the gas off in a \textbf{Flare Stack}.

\begin{center}
\small
\begin{tabular}{@{}lr@{}}
\toprule
Injectant & Recovered per mB \\
\midrule
Water                                & 0.35 \\
Sour gas, straight off the wellhead  & 0.55 \\
Natural gas, scrubbed                & 0.70 \\
\bottomrule
\end{tabular}
\end{center}

The middle row is the one that matters: a field can drive its own enhanced recovery on the waste
stream it was already producing. Recovery is capped per deposit --- by default a permanent
allowance of 20\,\% of the field's original capacity, which a server restart does not refund --- so
this extends a field rather than making it eternal.

\section{Refining}

Crude is separated into its fractions by the \textbf{Distillation Tower}, and a \textbf{Cracking
Unit} converts between them. The cracker has no temperature dial --- the dial is the firebox. Run
it hot on heavy fuel oil and it yields gasoline; run it cool on gas and it yields diesel. A knob
built out of blocks.

\subsection*{Industrial Burner}

A $3\times3\times3$ firebox of blast brick under a steel sheetmetal crown, with an Oilfield Frame
burner head in the front face where the fuel line lands and a second frame in the middle of the
crown for the flue. It makes \emph{no Flux at all}. What it makes is heat, and the crown is the
working face: whatever the burner serves sits on top of it.

Its reason to exist is heavy fuel oil, which is deliberately not a Diesel Generator fuel. The
burner is what makes owning a barrel of the stuff better than owning nothing.

\begin{ietip}
Build it whichever way round you like. Formation tries all four facings and takes the one you
actually built, so a correct burner is never refused for having been approached from the wrong
side.
\end{ietip}

\subsection*{Gas Scrubber}

A $3\times3\times6$ pair of vessels that sweetens sour gas: 100\,mB of sour gas yields 80\,mB of
natural gas, and every 2\,000\,mB scrubbed also drops a sulfur dust.

\textbf{It runs on heat, not Flux.} It draws heat from an Industrial Burner standing against its
skid, and does nothing whatever without one. That is also the prettiest loop in the chain --- a
burner will run on the scrubber's own natural gas, so the plant bootstraps itself; feed the burner
heavy fuel oil instead and you keep all of the gas. Two waste streams cancelling out.

An underfed scrubber runs slowly rather than stalling, and its overlay names which of four states
it is in: idle, scrubbing, no heat, or backed up. Its comparator reports the sour feed level,
because a feed backing up means it is not keeping pace with the field.

\subsection*{Lubrication Manifold}

A single block bolted onto a working machine that rebates a flat \textbf{8\,FE/t} of its running
cost while it has lubricant in it, drawing 2\,mB per 20-tick batch --- about eight minutes of
continuous running per bucket.

This is what lubricant is \emph{for}: every other cut off the column burns, and this is the one
that does not. The rebate is flat and absolute rather than proportional, so it stays a helpful
discount on a big machine instead of turning one net-positive.

\begin{ietip}
It greases nothing while its host is idle, so it will not quietly drain a bucket into a machine
that has no work queued.
\end{ietip}

\section{Storage}

Three \textbf{buried tanks} are invisible except for a \textbf{Tank Fill Cap} at the surface, which
carries the gauge, the comparator level and the pipe connection. Build the shell, cover the roof,
and hammer the cap.

\begin{center}
\small
\begin{tabular}{@{}llrrl@{}}
\toprule
Tier & Box & Blocks & Capacity & Metal \\
\midrule
Domestic Tank   & $2\times2\times2$ & 8   & 32\,000 mB     & Iron sheetmetal \\
Commercial Tank & $5\times5\times4$ & 82  & 512\,000 mB    & Steel sheetmetal \\
Bulk Depot      & $9\times9\times6$ & 290 & 4\,000\,000 mB & Steel sheetmetal \\
\bottomrule
\end{tabular}
\end{center}

Capacity per block rises with the tier, which is the reason to dig the bigger hole. All three are
\emph{hollow} --- you are lining an excavation, not filling one.

\begin{ietip}
The tank forms only once every roof block but the cap is buried. If the shape is right and only
the cover is missing it will say so, rather than refuse silently, because a tank that looks
completely correct is a miserable thing to debug.
\end{ietip}

The \textbf{Storage Tank} is the other kind: a hollow box of any rectangular shape from
$3\times3\times3$ to $16\times16\times16$, holding sixteen buckets per enclosed block. Build the
shell you want and strike any wall with a hammer. Pipe into any wall block, not just one fitting.

\begin{ietip}
Buried tanks are for supply you do not want to look at --- under a forecourt, under a yard. The
Storage Tank is for a tank farm you do.
\end{ietip}

For propane specifically there are three single-block vessels: the \textbf{Propane Cylinder} at four
buckets, the \textbf{Upright Propane Tank} at twelve, and the \textbf{Torpedo Propane Tank} at
twenty-four. They differ only in size and shape --- all three take propane and nothing else, all
three keep their contents when broken, and all three fill from a bucket or a pipe. Pick the one that
suits what you are building: a barbecue, a workshop wall, or the end of a garden.

\section{Burning it}

A ladder of power plants, built along one axis: \textbf{response against efficiency}. At the bottom
sits the \textbf{Portable Generator}, a single block for a build site.

\begin{description}
\item[Reciprocating Engine Bank] $4\times5\times5$ of cylinders burning diesel, biodiesel or
      natural gas. Instant response, mediocre efficiency, and the only plant that scales by
      building another one --- place a second bank alongside and they run as a set.
\item[Gas Turbine] The top of the gas-burning line, and the machine the rest of the chain is
      pointed at.
\item[Fuel Oil Boiler] A $5\times5\times7$ furnace making \emph{steam, not Flux}. Water and fuel in
      along the base course, steam out of the crown. Crude and diesel burn in it too, both worse
      trades than the heavy fuel oil it exists for.
\item[Heat Recovery Steam Generator] A $3\times5\times3$ bank of finned tubes that burns nothing at
      all. Butted flush against a Gas Turbine's exhaust end, it makes steam from heat that was
      going up the stack anyway.
\item[Steam Turbine Hall] $5\times9\times5$, the largest structure in the fork and the most
      efficient thing in the mod at turning fuel into Flux. Steam in at the condenser end, Flux out
      at the switchyard.
\end{description}

Steam is a real fluid here precisely so that the boiler house and the hall can be two buildings
plumbed together across a yard, rather than one enormous machine.

\subsection*{The combined cycle}

Gas Turbine $\rightarrow$ HRSG $\rightarrow$ Steam Turbine Hall. The turbine's waste heat becomes
steam and the steam becomes more Flux, so a millibucket of gas does considerably more work than it
does alone. It is also the least forgiving thing here to build: three structures in a line, one
fuel line, one steam line, and any of the three idle makes the other two worse.

\begin{ietip}
Only one HRSG may claim a given turbine's exhaust --- a second bolted to the same machine gets
nothing, because the waste heat is a fixed quantity and stacking would multiply it out of nowhere.
The HRSG is three by three because that is exactly the turbine's exhaust face: if it will not sit
flush, it is not in line.
\end{ietip}

\begin{iewarning}
The Steam Turbine Hall is deliberately slow to start and punished for cycling. Fed steadily it
beats everything; switched on and off against a wobbling demand it spends most of its life
spinning up.
\end{iewarning}

\section{Roads}

The bottom of the barrel is \textbf{bitumen}, and paving is what stops it being a waste product.
Put 250\,mB of bitumen through a \textbf{Mixer} with one sand and two gravel for 500\,mB of
\textbf{Wet Asphalt}, then pour it. It flows, finds its own level, and a moment later sets into a
full block of \textbf{Asphalt}.

There are no forms and no depth to judge, because a road is flat by definition: however sloppily it
is poured, it always sets flat. The mess is in the paving, not in the result. Set asphalt is a
little slippery --- enough that a road reads as faster than the mud beside it, without being ice.

Two further finishes come off the plain block. Four of them in a $2\times2$ give four
\textbf{Asphalt Tiles}, a laid pattern for forecourts and yards rather than open road; one of them
with a yellow dye gives \textbf{Asphalt with Markings}, carrying a painted lane stripe.

\section{The forecourt}

A \textbf{Gas Station Pump}, a \textbf{Forecourt Price Sign} and a nozzle that binds to a pump
within eight blocks. Fuel is dispensed to a person rather than to a machine.

\begin{iewarning}
Nothing past the Gas Turbine has been extensively playtested. The arithmetic is unit-tested and the
server starts clean, but the balance of the later chain is still unproven.
\end{iewarning}
